\section{Theorem-level relation to realization and coding}
\label{sec:literature29}
There are three different minimization problems in the argument, and
their classical precedents should not be conflated.

\paragraph{Positive realization versus an intrinsic complete class.}
Heller's positive-realization viewpoint and the HMM realization survey
of Vidyasagar \cite{Heller,Vidyasagar} concern the representation of
stochastic laws by positive state systems. They motivate the normal
form used here. Our main theorem does not replace that theory or
solve its unrestricted minimization problem. Its finite input is a
complete causal response array with declared cuts. From all its
normalized residuals, it computes a face-intersection graph and a sum
of affine ranks. Componentwise simpliciality is a finite geometric
test, and under that test the lower sum equals a simultaneously
realizable profile. No choice of a witness family remains in the
answer. The general normal-form statement alone does not supply this
all-realizations lower bound or its completeness test.

\paragraph{Residual minimality versus all-machine minimality.}
Denis and Esposito \cite{DenisEsposito}, Proposition 16 and their
probabilistic residual-automaton results, establish canonical minimal
generating subsets of residuals under their finite-generation
hypotheses. The residual-vertex upper construction in our proof belongs
to that line of ideas; its existence is not claimed as a new principle.
But being smallest among residual-state representations need not mean
being smallest among all positive representations. Our lower proof
allows arbitrary state continuations outside the residual hull and
rules out their sharing across distinct support components. It is
this additional step, together with affine independence inside each
component, that equates residual size with unrestricted stochastic
size in the stated class. The square example after
Corollary~\ref{cor:intrinsic-schedule29} exhibits failure of rank
completeness outside that class.

\paragraph{Nested polytopes versus causally compatible covers.}
Gillis and Glineur \cite{GG}, Definition 1 and Theorem 1, distinguish
restricted nonnegative rank and give its nested-polytope formulation.
Their distinction between ordinary and restricted rank is retained:
we never force a competing generator into the residual affine span.
For normalized binary-query channels the cube containment in
\eqref{eq:cube-factor29} follows directly by conditioning. That
checkpoint characterization is not a new identification of
nonnegative factorization with convex geometry. The additional online
question is whether covers at successive cuts admit stochastic shifts.
The explicit inclusions \eqref{eq:access-nesting29} answer it for every
preparation length and attain the lower rank at every acquisition cut.
A final enclosing simplex without those inclusions would not establish
Theorem~\ref{thm:access29}.

Reusch and Merzenich \cite{Reusch} study compatible coverings for
incompletely specified sequential machines. Their subject is particularly
relevant when a target has support zeros. We do not infer a canonical
completion from an almost-sure law. The present intrinsic theorem uses
a complete array, and the retained physical support-boundary proof
specifies its completion and supported tests explicitly.

\paragraph{Random-access codes versus exact online channel synthesis.}
Classical and quantum random-access coding, including the distinction
between private and shared randomness, is established in
\cite{ALMO}. Recent work studies average and worst-case coding
probabilities \cite{RAC2026,Liu2026}. The classical private-randomness
factorization in \cite{Liu2026}, equation (4), is precisely the
checkpoint conditioning mechanism underlying \eqref{eq:cube-factor29}.
Thus the random-access model and that factorization are not inventions
of this article. The present statement fixes every conditional
probability to \eqref{eq:access-law29} and prices the entire sequential
encoder. The affine-dimension lower bound, the compatible $(t+1)$-state
streaming construction, and the disjoint corner-cap converse give two
exact regimes with different memory scales in one growing full-support
family. We do not claim the optimal success probability for every
message size or a complete intermediate-signal phase diagram.

\paragraph{Statistical contact and the scope of originality.}
The new geometric classification applies to laws already specified or
forced. A minimax problem supplies those laws only after Bayes
optimality and uniform contact have both been imposed. The marked
application makes this distinction concrete: the cheaper three-event
Bayes tie rule is not uniformly minimax. Contact forces the two extra
coordinates before intrinsic realization is considered. This comparison
isolates the hypotheses and conclusions being used; it is not a claim
that no reformulation of earlier theory could imply them. No complete
priority or proof-level non-overlap audit of Norberg's filtered-experiment
comparison \cite{Norberg} is asserted. The finite geometric and online
results above are proved directly, with the inherited comparison,
control, and testing mechanisms attributed in the supporting text.
